% 彗星（综述）
% license CCBYSA3
% type Wiki

本文根据 CC-BY-SA 协议转载翻译自维基百科\href{https://en.wikipedia.org/wiki/Comet}{相关文章}。

\begin{figure}[ht]
\centering
\includegraphics[width=8cm]{./figures/45ca94c34375617b.png}
\caption{} \label{fig_Comet_1}
\end{figure}
彗星是一类由冰组成的、体积较小的太阳系天体或星际天体。当其经过接近太阳的区域时，会因受热而开始释放气体，这一过程称为释气作用（outgassing）。这一过程会在彗核周围产生一个延展的、未被引力束缚的大气层，即彗发（coma），并有时形成由彗发中吹离的气体和尘埃构成的彗尾。这些现象源于太阳辐射以及向外流动的太阳风等离子体对彗核的作用。彗核的大小从数百米到数十公里不等，由松散聚集的冰、尘埃和小型岩石颗粒组成。彗发的尺度可达地球直径的 15 倍，而彗尾的长度甚至可能超过一个天文单位。如果距离足够近且亮度足够高，人类在无望远镜的情况下即可从地球上观测到彗星，其在天空中的角距可达 30°（相当于 60 个月亮）。自古以来，彗星已被众多文化和宗教观察并记录。

彗星通常具有高离心率的椭圆轨道，其轨道周期范围极广，从数年到可能长达数百万年。短周期彗星来源于海王星轨道之外的柯伊伯带或相关的散射盘。长周期彗星则被认为起源于奥尔特云，一个由冰质天体组成的球形云层，其范围从柯伊伯带外侧延伸至通往最近恒星方向一半的距离\(^\text{[2]}\)。长周期彗星通常因经过恒星的引力扰动或银河潮汐效应而被触发向太阳方向运动。双曲轨道彗星可能只会一次穿越内太阳系，随后被抛入星际空间。彗星的出现称为一次显现（apparition）。

多次接近太阳的灭绝彗星会失去几乎全部挥发性冰与尘埃，并可能逐渐呈现出类似小型小行星的外观\(^\text{[3]}\)。一般认为，小行星与彗星的起源不同：小行星形成于木星轨道以内，而非太阳系外部\(^\text{[4][5]}\)。然而，主带彗星（main-belt comets）以及活跃的半人马族小天体的发现，使小行星与彗星之间的界限变得模糊。21 世纪初，一些具备长周期彗星轨道却呈现内太阳系小行星特征的小天体被称为曼岛彗星（Manx comets）。它们仍被归类为彗星，例如 C/2014 S3（PANSTARRS）\(^\text{[6]}\)。在 2013 年至 2017 年之间，共发现了 27 颗曼岛彗星\(^\text{[7]}\)。

截至 2021 年 11 月，共已知 4,584 颗彗星\(^\text{[8]}\)。然而，这仅占可能存在的彗星总量的一小部分，因为太阳系外部（奥尔特云）中类彗星天体的储库估计约为一万亿个\(^\text{[9][10]}\)。平均每年约有一颗彗星可被肉眼看见，尽管其中许多较为暗淡、外观并不特别醒目\(^\text{[11]}\)。特别明亮的彗星被称为“伟大彗星”（great comets）。人类已利用无人探测器访问过彗星，例如美国 NASA 的“深度撞击号”（Deep Impact），其通过撞击坦普尔 1 号彗星形成陨坑以研究其内部结构；又如欧洲航天局的“罗塞塔号”（Rosetta），它成为首个在彗星上着陆的机器人探测器\(^\text{[12]}\)。
\subsection{词源}
\begin{figure}[ht]
\centering
\includegraphics[width=8cm]{./figures/6a3173d4cf8911ea.png}
\caption{公元 729 年的《盎格鲁－撒克逊编年史》以及尊敬的比德（the Venerable Bede）的著作中均提到了彗星。} \label{fig_Comet_2}
\end{figure}
单词“comet（彗星）”源自古英语 cometa，其来源为拉丁语 comēta 或 comētēs，而这些拉丁词又是对希腊语 κομήτης 的罗马化形式，意为“留着长发的（事物）”。《牛津英语词典》指出，术语（ἀστὴρ） κομήτης 在希腊语中早已表示“长发之星，即彗星”。κομήτης 源自动词 κομᾶν（koman），意为“留长发”，而该动词又派生自名词 κόμη（komē），意为“头发”，并被用来表示“彗尾”\(^\text{[15][16]}\)。

彗星的天文学符号（在 Unicode 中表示）为 U+2604 ☄ COMET，由一个小圆盘及其后三条类似发丝的延伸部分组成\(^\text{[17]}\)。
\subsection{物理特征}
\begin{figure}[ht]
\centering
\includegraphics[width=8cm]{./figures/37fee7dd52548b31.png}
\caption{彗星的结构} \label{fig_Comet_3}
\end{figure}
\subsubsection{彗核}
\begin{figure}[ht]
\centering
\includegraphics[width=8cm]{./figures/43510b130873771c.png}
\caption{在航天器飞越期间拍摄的 103P/Hartley 彗核图像。该彗核长度约为 2 千米。} \label{fig_Comet_4}
\end{figure}
彗星的固体核心结构被称为彗核。彗核由岩石、尘埃、水冰以及冻结的二氧化碳、一氧化碳、甲烷和氨等物质的混合体构成\(^\text{[18]}\)。因此，根据弗雷德·惠普尔（Fred Whipple）的模型，它们常被形象地称为“肮脏的雪球”（dirty snowballs）\(^\text{[19]}\)。尘埃含量更高的彗星则被称为“冰质尘球”（icy dirtballs）\(^\text{[20]}\)。“冰质尘球”这一术语源于 2005 年 7 月 NASA“深度撞击号”（Deep Impact）任务向坦普尔 1 号彗星（Comet 9P/Tempel 1）发送“撞击器”探测器并观测其碰撞后的结果。2014 年的一项研究表明，彗星类似于“炸冰淇淋”，即其表面由密集的晶体冰与有机化合物混合构成，而内部冰层则更冷且密度更低\(^\text{[21]}\)。

彗核表面通常干燥、多尘或岩质，表明冰层隐藏在厚度数米的表层壳体之下。彗核中含有多种有机化合物，其中可能包括甲醇、氰化氢、甲醛、乙醇、乙烷，以及可能更复杂的分子，如长链烃类与氨基酸\(^\text{[22][23]}\)。2009 年确认，NASA “星尘号”（Stardust）任务回收的彗星尘埃中发现了氨基酸甘氨酸\(^\text{[24]}\)。2011 年 8 月，一份基于 NASA 对地球陨石研究的报告提出，DNA 与 RNA 的组成单元（腺嘌呤、鸟嘌呤及相关有机分子）可能形成于小行星和彗星\(^\text{[25][26]}\)。

彗核外层表面具有极低反照率，使其成为太阳系中反射率最低的天体之一。贾托号（Giotto）探测器发现，哈雷彗星（1P/Halley）的彗核仅反射约四个百分点的入射光\(^\text{[27]}\)；而“深空一号”（Deep Space 1）发现波罗雷彗星（Comet Borrelly）的表面反射率甚至不足 3.0\%\(^\text{[27]}\)。相比之下，沥青可反射约 7\%。彗核表面的暗色物质可能由复杂的有机化合物构成。太阳加热会驱散较轻、易挥发的化合物，而留下更大的有机分子，这些物质通常极为黝黑，如焦油或原油。彗核表面的低反照率使其能够吸收驱动释气过程的热量\(^\text{[28]}\)。

已观测到的彗核半径可达 30 千米（19 英里）\(^\text{[29]}\)，但要准确测定彗核的大小并不容易\(^\text{[30]}\)。322P/SOHO 的彗核直径可能只有 100–200 米（330–660 英尺）\(^\text{[31]}\)。尽管观测仪器的灵敏度不断提高，但仍缺乏对更小彗星的探测，这使部分研究者推测直径小于 100 米（330 英尺）的彗星数量可能确实稀少\(^\text{[32]}\)。已知彗星的平均密度估计约为 $0.6 g/cm^3$（0.35 oz/cu in）\(^\text{[33]}\)。由于质量较低，彗核不会在自引力作用下成为球形，因此呈现不规则形状\(^\text{[34]}\)。
\begin{figure}[ht]
\centering
\includegraphics[width=8cm]{./figures/7e9645911c29b572.png}
\caption{彗星 81P/Wild 在光照面和阴暗面都呈现出喷流，其表面起伏鲜明，且十分干燥。} \label{fig_Comet_5}
\end{figure}
大约六个百分点的近地小行星被认为是不再发生释气作用的彗星灭绝彗核\(^\text{[35]}\)，其中包括 14827 Hypnos 与 3552 Don Quixote。

“罗塞塔号”（Rosetta）与“菲莱”（Philae）着陆器的探测结果显示，67P/丘留莫夫－格拉西缅科彗星（67P/Churyumov–Gerasimenko）的彗核不存在磁场\(^\text{[36][37]}\)，这表明磁性可能并未在早期微行星形成过程中发挥作用。此外，罗塞塔号上的 ALICE 光谱仪确定，在距离彗核约 1 千米范围内，由太阳辐射对水分子进行光电离所产生的电子，而非此前认为的来自太阳的光子，是导致彗核释放出的水与二氧化碳分子在彗发中发生分解的主要因素\(^\text{[38][39]}\)。菲莱着陆器上的仪器在彗星表面检测到至少十六种有机化合物，其中有四种（乙酰胺、丙酮、异氰酸甲酯和丙醛）是首次在彗星上被发现的\(^\text{[40][41][42]}\)。
\begin{figure}[ht]
\centering
\includegraphics[width=10cm]{./figures/8c2117a476d375ac.png}
\caption{} \label{fig_Comet_6}
\end{figure}
\subsubsection{彗发}
\begin{figure}[ht]
\centering
\includegraphics[width=8cm]{./figures/3f4fb64d5886d128.png}
\caption{Borrelly 彗星呈现喷流，但其表面不存在冰。} \label{fig_Comet_7}
\end{figure}
\begin{figure}[ht]
\centering
\includegraphics[width=8cm]{./figures/2d693476f2b0c944.png}
\caption{近日点前不久由哈勃望远镜拍摄的 ISON 彗星图像[50]} \label{fig_Comet_8}
\end{figure}
由此释放出的尘埃和气体会在彗星周围形成一个巨大而极其稀薄的大气层，称为“彗发”（coma）。太阳的辐射压与太阳风施加在彗发上的力会形成一条巨大的彗尾，且彗尾始终指向远离太阳的方向\(^\text{[51]}\)。

彗发通常由水和尘埃组成，其中水占彗星在距太阳 3 至 4 个天文单位（4.5 亿至 6 亿千米；2.8 亿至 3.7 亿英里）范围内从彗核中流出的挥发物的 90\% 之多\(^\text{[52]}\)。母体分子 $H_2O$主要通过光解作用被破坏，而光电离的贡献要小得多；与光化学过程相比，太阳风对水的破坏作用仅占很小比例\(^\text{[52]}\)。较大的尘埃颗粒会沿着彗星轨道路径残留，而较小的颗粒则被阳光压力推向远离太阳的方向进入彗尾\(^\text{[53]}\)。

虽然彗星的固体彗核通常小于 60 千米（37 英里）宽，但彗发的大小可能达到数千甚至数百万千米，有时甚至比太阳还大\(^\text{[54]}\)。例如，2007 年 10 月的一次爆发过后大约一个月，彗星 17P/Holmes 的稀薄尘埃大气层短暂地变得比太阳还大\(^\text{[55]}\)。1811 年的“大彗星”其彗发的直径大约与太阳相当\(^\text{[56]}\)。尽管彗发可以变得非常庞大，但其尺寸会在彗星经过火星轨道附近（约 1.5 AU，即 2.2 亿千米；1.4 亿英里）时缩小\(^\text{[56]}\)。在这一距离，太阳风强度足以将彗发中的气体和尘埃吹散，从而反而使彗尾变得更长\(^\text{[56]}\)。离子彗尾的长度可达到一个天文单位（1.5 亿千米）或更长\(^\text{[55]}\)。
\begin{figure}[ht]
\centering
\includegraphics[width=8cm]{./figures/f1b7d1bdbd91578d.png}
\caption{C/2006 W3（Christensen）释放碳气体（红外图像）} \label{fig_Comet_9}
\end{figure}
彗发与彗尾都由太阳照亮，并可能在彗星穿过内太阳系时变得可见；其中尘埃直接反射阳光，而气体则因电离而产生辉光\(^\text{[57]}\)。大多数彗星过于暗淡，以至于没有望远镜无法看见，但每隔数十年会有少数彗星亮到可被肉眼观测\(^\text{[58]}\)。有时彗星会经历一次巨大而突然的气体与尘埃爆发，在此期间彗发的尺寸会显著增大，持续一段时间。2007 年的霍姆斯彗星（Comet Holmes）便发生了这一现象\(^\text{[59]}\)。

1996 年，人们发现彗星能够发射 X 射线\(^\text{[60]}\)。这一发现令天文学家非常惊讶，因为 X 射线通常与极高温的天体相关。Thomas E. Cravens 在 1997 年初首次提出了解释\(^\text{[61]}\)。X 射线的产生源于彗星与太阳风之间的相互作用：当高电荷态的太阳风离子穿过彗发大气层时，它们会与彗星的原子和分子发生碰撞，并在这一过程中“夺取”一个或多个电子，这一机制称为“电荷交换”（charge exchange）。电子被太阳风离子俘获之后，会通过向基态去激发的过程发射出 X 射线和远紫外光子\(^\text{[62]}\)。
\subsubsection{弓形激波}
弓形激波是由太阳风与彗星电离层相互作用形成的，而彗星的电离层则由彗发中气体的电离所产生。随着彗星逐渐靠近太阳，释气速率不断增加，使得彗发膨胀；阳光会进一步电离彗发中的气体。当太阳风经过这一离子化的彗发区域时，弓形激波便会出现。

第一次观测是在 1980 年代与 1990 年代，当时多艘航天器先后飞掠了 21P/Giacobini–Zinner\(^\text{[63]}\)、1P/Halley\(^\text{[64]}\)与 26P/Grigg–Skjellerup\(^\text{[65]}\)彗星。结果发现，彗星的弓形激波比例如地球那样的行星弓形激波更宽且更加平缓。这些观测均发生在近日点附近，当时弓形激波已完全形成。

罗塞塔号（Rosetta）在 67P/Churyumov–Gerasimenko 彗星上观测到弓形激波的早期形成阶段，当时彗星正向太阳靠近，释气活动逐渐增强。这个尚在成长中的弓形激波被称为“婴儿弓形激波”（infant bow shock）。婴儿弓形激波是不对称的，并且相对于彗核距离而言，比完全发育的弓形激波更为宽广\(^\text{[66]}\)。
\subsubsection{彗尾}
\begin{figure}[ht]
\centering
\includegraphics[width=8cm]{./figures/ec7df3f317631c33.png}
\caption{彗星在接近太阳轨道时彗尾的典型指向方向} \label{fig_Comet_10}
\end{figure}
在外太阳系中，彗星保持冻结且处于不活跃状态，由于体积很小，从地球上极难或几乎不可能被探测到。基于哈勃空间望远镜的观测，有研究报告在柯伊伯带中统计性地探测到不活跃的彗星彗核[67][68]，但这些探测结果曾受到质疑[69][70]。当彗星进入内太阳系时，太阳辐射会使其中的挥发性物质升华并从彗核中逸出，同时携带尘埃。

逸出的尘埃与气体分别形成各自独立的彗尾，方向略有不同。尘埃彗尾通常沿彗星轨道附近被遗留，从而形成一条常常呈弯曲状的彗尾，称为 II 型或尘埃彗尾（dust tail）[57]。同时，离子彗尾（ion tail），亦称 I 型彗尾，由气体组成，始终直接指向远离太阳的方向，因为气体比尘埃更强烈地受到太阳风作用，并沿磁力线运动，而非遵循轨道轨迹[71]。在某些情况下——例如当地球经过彗星的轨道平面时——可以看到一条反尾（antitail），其方向与离子彗尾和尘埃彗尾相反[72]。
\begin{figure}[ht]
\centering
\includegraphics[width=8cm]{./figures/1dacb212df9b211e.png}
\caption{彗星示意图：展示由太阳风形成的尘埃轨迹、尘埃彗尾与离子气体彗尾} \label{fig_Comet_11}
\end{figure}
对反尾（antitail）的观测对太阳风的发现起到了重要作用[73]。离子彗尾是由于彗发中粒子被太阳紫外辐射电离而形成的。一旦粒子被电离，它们便获得净正电荷，从而在彗星周围产生一个“感应磁层”（induced magnetosphere）。彗星及其感应磁场对向外流动的太阳风粒子构成阻碍。由于彗星与太阳风之间的相对轨道速度为超音速，彗星在太阳风流向的上游会形成一个弓形激波。在这一弓形激波中，彗星离子（称为“拾取离子” pick-up ions）大量聚集，并使太阳磁场被等离子体“负载”，使磁力线在彗星周围“披覆”（drape），从而形成离子彗尾[74]。

若离子彗尾的负载程度足够大，磁力线会被压缩到某一距离处发生磁重联（magnetic reconnection）。这会导致一次“彗尾断裂事件”（tail disconnection event）[74]。这一现象已多次被观测到，其中一个著名事件发生在 2007 年 4 月 20 日：当恩克彗星（Encke’s Comet）穿过一次日冕物质抛射（CME）时，其离子彗尾完全被切断。该事件由 STEREO 空间探测器记录[75]。

2013 年，欧洲航天局（ESA）的科学家报告称，金星的电离层会向外流出，其行为与类似条件下彗星的离子彗尾相似[76][77]。
\subsubsection{喷发}
\begin{figure}[ht]
\centering
\includegraphics[width=8cm]{./figures/e97b94991a13b0fc.png}
\caption{103P/Hartley 的气体与冰雪喷流} \label{fig_Comet_12}
\end{figure}
不均匀的加热会使新生成的气体从彗核表面的薄弱区域喷出，类似于间歇泉[78]。这些气体和尘埃的喷流能够使彗核旋转，甚至将其分裂[78]。2010 年的研究显示，干冰（冻结的二氧化碳）的升华能够为从彗核喷出的物质喷流提供动力[79]。对 Hartley 2 彗星的红外成像显示了此类喷流从其表面喷出，并携带尘埃颗粒进入彗发[80]。
\subsection{轨道特征}
大多数彗星是太阳系内的小天体，其轨道为细长的椭圆轨道，轨道的一部分将它们带到靠近太阳的区域，而其余时间则位于太阳系的遥远外侧[81]。彗星通常根据其轨道周期长短进行分类：周期越长，其轨道椭圆的离心率越大。









